\section{\texorpdfstring{Karnaughova mapa}{Karnaughova mapa}}

% =====================================================================================================================
	\begin{frame}
    \frametitle{Návrh mapy}
		\boldmath
		
		\begin{minipage}[t][\textheight][t]{0.48\textwidth}
		\vspace{5mm}
		\small
		\textbf{Základní buňka - 2 proměnné:}

		\begin{karnaugh-map}[2][2][1][$B$][$A$]

			% ====== vlastní čáry ======

			% A = 1 (druhý řádek)
			\draw[very thick] (-0.15,1) -- (-0.15,0);

			% B = 1 (sloupce 3 a 4)
			\draw[very thick] (1,2.15) -- (2,2.15);

		\end{karnaugh-map}
      
		\end{minipage}
		\begin{minipage}[t][\textheight][t]{0.45\textwidth}
    \vspace{5mm}
		\small
		\textbf{Zrcadlení - 3 proměnné:}
		
		\begin{karnaugh-map}[4][2][1][$B$][$C$][$A$]

			% ====== vlastní čáry ======

			% A = 1 (druhý řádek)
			\draw[very thick] (-0.15,1) -- (-0.15,0);

			% B = 1 (sloupce 3 a 4)
			\draw[very thick] (2,2.3) -- (4,2.3);

			% C = 1 (sloupce 2 a 3)
			\draw[very thick] (1,2.15) -- (3,2.15);

		\end{karnaugh-map}
    Nad přidanými buňkami je stav nové proměnné roven jedné.
		\end{minipage}
		
	\end{frame}
	
% =====================================================================================================================
	\begin{frame}
    \frametitle{Návrh mapy - 4 proměnné}
		\boldmath
		
		\small
		\textbf{Poslední zrcadlení je směrem dolů - 4 proměnné:}
		
		
		\begin{minipage}[t][0.8\textheight][t]{0.5\textwidth}
		\vspace{3mm}
			\begin{karnaugh-map}[4][4][1][$B$][$C$][$A$][$D$]

				% ====== vlastní čáry ======

				% A = 1 (řádek 2 a 3)
				\draw[very thick] (-0.15,1) -- (-0.15,3);

				% B = 1 (sloupce 3 a 4)
				\draw[very thick] (2,4.3) -- (4,4.3);

				% C = 1 (sloupce 2 a 3)
				\draw[very thick] (1,4.15) -- (3,4.15);
				
				% A = 1 (řádek 3 a 4)
				\draw[very thick] (-0.30,0) -- (-0.30,2);

			\end{karnaugh-map}
		\end{minipage}
		\begin{minipage}[t][0.8\textheight][t]{0.45\textwidth}
			
			\begin{itemize}
				\item Stejným způsobem vytváříme Grayův kód.
				\item Postup lze opakovat pro více proměnných, vzniká ale problém se sousedstvím jedniček v nesousedních částech mapy.
			\end{itemize}
		\end{minipage}
	\end{frame}
	
% =====================================================================================================================
	\begin{frame}
    \frametitle{Sousední jedničky uvnitř mapy}
		\boldmath
		\begin{center}
		  \begin{karnaugh-map}[4][4][1][$B$][$C$][$A$][$D$]
				\minterms{4,5,7,6,15}
				\implicant{4}{6}
				\implicant{7}{15}
				
				% ====== vlastní čáry ======

				% A = 1 (řádek 2 a 3)
				\draw[very thick] (-0.15,1) -- (-0.15,3);

				% B = 1 (sloupce 3 a 4)
				\draw[very thick] (2,4.3) -- (4,4.3);

				% C = 1 (sloupce 2 a 3)
				\draw[very thick] (1,4.15) -- (3,4.15);
				
				% A = 1 (řádek 3 a 4)
				\draw[very thick] (-0.30,0) -- (-0.30,2);

			\end{karnaugh-map}
			
			$Y = \textcolor{red}{A\bar{D}} + \textcolor{green}{ABC}$
		\end{center}
	\end{frame}
	
% =====================================================================================================================
	\begin{frame}
    \frametitle{Sousední jedničky přes okraj mapy}
		\boldmath
		\begin{center}
		  \begin{karnaugh-map}[4][4][1][$B$][$C$][$A$][$D$]
				\minterms{4,5,6,12,13}
				\implicant{4}{13}
				\implicantedge{4}{4}{6}{6}
				
				% ====== vlastní čáry ======

				% A = 1 (řádek 2 a 3)
				\draw[very thick] (-0.15,1) -- (-0.15,3);

				% B = 1 (sloupce 3 a 4)
				\draw[very thick] (2,4.3) -- (4,4.3);

				% C = 1 (sloupce 2 a 3)
				\draw[very thick] (1,4.15) -- (3,4.15);
				
				% A = 1 (řádek 3 a 4)
				\draw[very thick] (-0.30,0) -- (-0.30,2);

			\end{karnaugh-map}
			
			$Y = \textcolor{red}{A\bar{C}} + \textcolor{green}{A\bar{B}\bar{D}}$
		\end{center}
	\end{frame}

% =====================================================================================================================
	\begin{frame}
    \frametitle{Sousední jedničky v rozích mapy}
		\boldmath
		\begin{center}
		  \begin{karnaugh-map}[4][4][1][$B$][$C$][$A$][$D$]
				\minterms{0,2,8,10}
				\implicantcorner
				
				% ====== vlastní čáry ======

				% A = 1 (řádek 2 a 3)
				\draw[very thick] (-0.15,1) -- (-0.15,3);

				% B = 1 (sloupce 3 a 4)
				\draw[very thick] (2,4.3) -- (4,4.3);

				% C = 1 (sloupce 2 a 3)
				\draw[very thick] (1,4.15) -- (3,4.15);
				
				% A = 1 (řádek 3 a 4)
				\draw[very thick] (-0.30,0) -- (-0.30,2);

			\end{karnaugh-map}
			
			$Y = \textcolor{red}{\bar{A}\bar{B}}$
		\end{center}
	\end{frame}
	
% =====================================================================================================================
	\begin{frame}
    \frametitle{Nadbytečná skupina}
		\boldmath
		\begin{center}
		  \begin{karnaugh-map}[4][4][1][$B$][$C$][$A$][$D$]
				\minterms{4,5,13,15}
				\implicant{4}{5}
				\implicant{13}{15}
				\implicant{5}{13}
				
				% ====== vlastní čáry ======

				% A = 1 (řádek 2 a 3)
				\draw[very thick] (-0.15,1) -- (-0.15,3);

				% B = 1 (sloupce 3 a 4)
				\draw[very thick] (2,4.3) -- (4,4.3);

				% C = 1 (sloupce 2 a 3)
				\draw[very thick] (1,4.15) -- (3,4.15);
				
				% A = 1 (řádek 3 a 4)
				\draw[very thick] (-0.30,0) -- (-0.30,2);

			\end{karnaugh-map}
			
			Minimalizace jen pro: $Y = \textcolor{red}{A\bar{C}\bar{D}} + \textcolor{green}{ABD}$
		\end{center}
	\end{frame}